\documentclass[10pt,a4paper]{amsart}
\usepackage[margin=19mm]{geometry}
\usepackage{amsmath,amssymb,mathtools}
\usepackage[scr=rsfs,cal=euler]{mathalfa}
\usepackage{xfrac}
\usepackage[unicode,hidelinks,bookmarks=false]{hyperref}
\newcommand{\ie}{i.e.\ }
\newcommand{\cf}{cf.\ }
\newcommand{\resp}{respectively\ }
\newcommand{\Subsection}{\subsection}
\providecommand{\nobreakdash}{\nobreak\hbox{-}}
\setlength{\emergencystretch}{2em}
\multlinegap0pt
\usepackage{amssymb}
\newtheorem{lemma}{Lemma}
\newcommand{\dif}{\mathrm{d}}
\newcommand{\majorarc}{\mathfrak{M}}
\title{Prime pairs along rays of prime indices}
\author{Shenghao Hua, Sizhe Xie}
\date{Source: arXiv:2609.07371v1 (CC BY 4.0)}
\begin{document}
\maketitle

\noindent\emph{Source and licence.} Prime pairs along rays of prime indices. Version arXiv:2609.07371v1 (CC BY 4.0), distributed by arXiv under the Creative Commons Attribution 4.0 International licence, http://creativecommons.org/licenses/by/4.0/, as recorded in the arXiv licence metadata for this exact version and shown on the abstract page https://arxiv.org/abs/2609.07371v1. Full licence notice: \texttt{licenses/arxiv-2609.07371-CC-BY-4.0.txt}.\par\smallskip

\noindent\textit{Editorial scope.} Counted unit: the article's lemma lem:principal-arcs (source0.tex lines 988--1008). The article's complete original proof of it is delivered with it (source0.tex lines 1010--1165, one \texttt{proof} environment of 156 lines). No mathematics is written, repaired or re-derived: the statement and the proof are the article's own lines, in the article's own notation, and no retained line is substituted or re-flowed.  Retained uncounted as the source-local context the proof needs: lines 505--507; lines 551--554; lines 790--806.  Nothing was omitted from any retained span, so the package carries no omission record.  No retained line contains a citation, so the wrapper carries no bibliography and no bibliography entry.  No retained line contains a figure environment, an image-inclusion command or drawing code, so no image is delivered and the package declares no asset dependency.  Editorial formatting only, in addition: an AMS \texttt{amsart} preamble that restates the article's own declarations and macros used by the retained lines, namely the environments \texttt{lemma} on separate counters, together with the standard packages those lines need.  The retained environments print in this document as Lemma 1 in order of appearance, whereas the article prints the same items with the numbering of its own full document.  The complete native source of the article as distributed in the arXiv e-print for this exact version is bundled unchanged as \texttt{sources/209622/source0.tex}.
\begin{equation}\label{three-cutoffs}
 z_+=x^2+y,\qquad z_-=x^2,\qquad z_0=y.
\end{equation}

\begin{equation}\label{circle-parameters}
 Q_1=(\log x)^c,\qquad Q=x^{1-\delta},
 \qquad 0<\delta<\frac1{10},
\end{equation}

\begin{lemma}[Singular series tails]\label{lem:tail-bound}
Let $\varkappa\in\{-1,+1\}$ and
\[
 \mathcal T_\varkappa(k,Q_1)=
 \lim_{V\to\infty}
 \sum_{\substack{Q_1<q\leq V\\2\nmid q}}
 \frac{\mu(q)}{\varphi(q)}
 \left(\frac{\varkappa k}{q}\right).
\]
Recall that $Q_1=(\log x)^c$ in \eqref{circle-parameters}.
For every $C>0$, we can take $c=c(C,A,K_0)$ sufficiently large.  Then
\begin{equation}\label{tail-second-moment}
 \sum_{\substack{k\leq y\\\mu^2(k)=1\\k>1}}
 |\mathcal T_\varkappa(k,Q_1)|^2
 \ll_{C,A,K_0}\frac{y}{(\log x)^C}.
\end{equation}
\end{lemma}

\begin{lemma}[Principal major arcs]\label{lem:principal-arcs}
For every $C>0$, provided that $c$ in $Q_1=(\log x)^c$
is sufficiently large in terms of $C$, $A$, and $K_0$, one has
\begin{equation}\label{principal-L2}
 \sum_{\substack{k\leq y\\\mu^2(k)=1\\k>1}}
 \left|\int_{\majorarc}T_{1,z_\nu}(\alpha)T_{2,s_\nu}(\alpha)
 e(t_\nu k\alpha)\,\dif\alpha
 -\mathfrak S_\nu(k)H_\nu(x,k)\right|^2
 \ll_{C,A,K_0}
 \frac{yx^2}{(\log x)^C}+yxQ(\log x)^{C_1},
\end{equation}
where
\[
 \mathfrak S_\nu(k)=
 \begin{cases}
  \mathfrak S_+(k),&\nu=+,\\
  \mathfrak S_-(k),&\nu=-,0,
 \end{cases}
\]
and where $C_1$ depends only on the fixed circle method parameters.
\end{lemma}

\begin{proof}
Fix $q$, $a$, and $d$.  The $\beta$ integral in the principal product
is
\[
 \int_{|\beta|\leq1/(qQ)}
 \left(\sum_{m\leq z_\nu}e(\beta m)\right)
 \left(\sum_{\substack{h\leq x\\(h,q)=d}}e(s_\nu\beta h^2)\right)
 e(t_\nu k\beta)\,\dif\beta.
\]
We extend this integral to a full period.  Orthogonality then gives
$m+s_\nu h^2+t_\nu k=0$.  The upper condition $m\leq z_\nu$ follows
from \eqref{three-cutoffs}.  The condition $m\geq1$ gives the intervals
\[
 \mathcal I_+(k)=[1,x],\qquad
 \mathcal I_-(k)=(\sqrt k,x],\qquad
 \mathcal I_0(k)=[1,x]\cap[1,\sqrt{k}).
\]
M\"obius inversion gives the following estimate uniformly for
$d\mid q$.
\begin{equation}\label{unified-coprime-count}
 \#\{h\in\mathcal I_\nu(k):(h,q)=d\}
 =\frac{\varphi(q/d)}qH_\nu(x,k)+O(\tau(q/d)).
\end{equation}

The moving intervals appear only after orthogonality on the full
period.  For the extension error, the quadratic sum is still the full
sum over $h\leq x$.  The factor $e(t_\nu k\beta)$ has absolute value
one.  A change of the sign $s_\nu$ only conjugates the quadratic sum.
Thus the estimate is uniform in $\nu$ and $k$.  Parseval's identity gives
\begin{equation}\label{quadratic-parseval}
 \int_0^1\left|
 \sum_{\substack{h\leq x\\(h,q)=d}}e(s\beta h^2)
 \right|^2\,\dif\beta
 =\#\{h\leq x:(h,q)=d\}\leq\frac xd+1.
\end{equation}
Outside the short arc, we have
$|\sum_{m\leq z}e(\beta m)|\ll\|\beta\|^{-1}$.  Cauchy's inequality
and \eqref{quadratic-parseval} give
\begin{equation}\label{beta-extension}
 \int_{1/(qQ)\leq|\beta|\leq1/2}
 \frac1{\|\beta\|}
 \left|\sum_{\substack{h\leq x\\(h,q)=d}}e(s_\nu\beta h^2)\right|
 \,\dif\beta
\ll(qQ)^{1/2}\left(\frac xd+1\right)^{1/2}.
\end{equation}
We have
\[
 |G_s(a,q,d)|\leq\varphi(q^*_{d}).
\]
Thus the total contribution of the error in
\eqref{unified-coprime-count} is
\begin{align*}
 \mathcal E_{\mathrm{count}}
 &\leq \sum_{q\leq Q_1}\frac{|\mu(q)|}{\varphi(q)}
       \sum_{a\bmod q}^{*}\sum_{d\mid q}
       \frac{|G_{s_\nu}(a,q,d)|}{\varphi(q^*_{d})}
       \tau(q/d)\notag\\
 &\ll \sum_{q\leq Q_1}\sum_{d\mid q}\tau(q/d)
 \ll Q_1(\log(2Q_1))^2
 \ll(\log x)^{C_1}.                            \label{counting-error-sum}
\end{align*}
The sum over $a$ cancels the factor $1/\varphi(q)$.  Moreover,
\[
 \frac{|G_{s_\nu}(a,q,d)|}{\varphi(q^*_{d})}\leq1.
\]
In the same way,
\eqref{beta-extension} gives the following bound after summation over
$a$, $d$, and $q$.
\begin{align}
 \mathcal E_{\mathrm{ext}}
 &\ll Q^{1/2}\sum_{q\leq Q_1}q^{1/2}
       \sum_{d\mid q}\left(\frac xd+1\right)^{1/2}\notag\\
 &\ll (xQ)^{1/2}Q_1^{3/2}(\log(2Q_1))^2
 \ll (xQ)^{1/2}(\log x)^{C_1}.                 \label{extension-error-sum}
\end{align}
Here we used $Q_1=(\log x)^c$.  We enlarged $C_1$ when needed.

We now find the main coefficient.
The factor $\mu(q)$ restricts $q$ to be squarefree.  Hence
$(d,q^*)=1$ for every $d\mid q$, and therefore
\[
 d_{q^*}=d,\qquad q^*_{d}=q^*=\frac qd.
\]
The classes with $(r,q)=d$
give
\[
 \sum_{d\mid q}G_s(a,q,d)
 =\sum_{r\bmod q}e\left(\frac{asr^2}{q}\right).
\]
Thus the main contribution for a fixed modulus $q$ is
\begin{equation*}\label{fixed-q-principal-contribution}
 \frac{\mu(q)H_\nu(x,k)}{q\varphi(q)}
 \sum_{a\bmod q}^{*}\sum_{r\bmod q}
 e\left(\frac{a(s_\nu r^2+t_\nu k)}q\right).
\end{equation*}
Define the rational complete sum as follows.
\[
 \Sigma_\nu(q,k)=
 \sum_{a\bmod q}^{*}\sum_{r\bmod q}
 e\left(\frac{a(s_\nu r^2+t_\nu k)}q\right).
\]
This sum is multiplicative.  At an odd prime $p$, we have
\[
 \Sigma_+(p,k)=p\left(\frac{-k}{p}\right),\qquad
 \Sigma_-(p,k)=\Sigma_0(p,k)=p\left(\frac{k}{p}\right),
\]
and all three sums vanish at $p=2$.  Hence, for odd squarefree $q$, we
have
\[
 \Sigma_\nu(q,k)=q\left(\frac{\varkappa_\nu k}{q}\right).
\]
We have proved the following formula for each $k$.
\begin{align}
 &\int_{\majorarc}T_{1,z_\nu}(\alpha)T_{2,s_\nu}(\alpha)
 e(t_\nu k\alpha)\,\dif\alpha\notag\\
 &\quad=H_\nu(x,k)
 \sum_{\substack{q\leq Q_1\\2\nmid q}}
 \frac{\mu(q)}{\varphi(q)}
 \left(\frac{\varkappa_\nu k}{q}\right)
 +O(\mathcal E_{\mathrm{count}}+\mathcal E_{\mathrm{ext}})
 \notag\\
 &\quad=H_\nu(x,k)
 \sum_{\substack{q\leq Q_1\\2\nmid q}}
 \frac{\mu(q)}{\varphi(q)}
 \left(\frac{\varkappa_\nu k}{q}\right)
 +O((xQ)^{1/2}(\log x)^{C_1}).                 \label{principal-pointwise}
\end{align}
Let $D_\nu(k)$ be the expression inside the absolute value on the left
of \eqref{principal-L2}.
By \eqref{principal-pointwise} and the definition of
$\mathcal T_{\varkappa_\nu}(k,Q_1)$ in Lemma \ref{lem:tail-bound},
\[
 D_\nu(k)
 =-H_\nu(x,k)\mathcal T_{\varkappa_\nu}(k,Q_1)
 +O(\mathcal E_{\mathrm{count}}+\mathcal E_{\mathrm{ext}}).
\]
Lemma~\ref{lem:tail-bound} and $H_\nu(x,k)\leq x$ now give
\begin{align*}
 \sum_{\substack{k\leq y\\\mu^2(k)=1\\k>1}}|D_\nu(k)|^2
 &\ll x^2\sum_{\substack{k\leq y\\\mu^2(k)=1\\k>1}}
       |\mathcal T_{\varkappa_\nu}(k,Q_1)|^2
       +y\mathcal E_{\mathrm{count}}^2
       +y\mathcal E_{\mathrm{ext}}^2\\
 &\ll \frac{yx^2}{(\log x)^C}
       +yxQ(\log x)^{C_1}.
\end{align*}
We enlarge $C_1$ once more.  It then absorbs the square of the
logarithmic factor in \eqref{extension-error-sum}.  Also, for every
fixed $C$ and all sufficiently large $x$, we have
\[
 y\mathcal E_{\mathrm{count}}^2
 \ll y(\log x)^{2C_1}
 \ll\frac{yx^2}{(\log x)^C}.
\]
This proves \eqref{principal-L2}.
\end{proof}

\end{document}
